\section{Overlap arrays and cascade calculus}\label{sec:arrays}

The rotation and cavity arguments use limits of several overlap arrays at
once. This section shows that the Ghirlanda--Guerra identities reduce these
limits to ordered scalar quantiles, and gives the Gaussian integration rules
for their finite-step approximations. We include proofs of the structural
facts needed below; they are instances of the ultrametricity and
synchronization theory developed in
\cite{PanchenkoUltra,PanchenkoMulti}. The cascade construction and its marking
calculus are closely related to \cite{BolthausenSznitman}.

\subsection{The Ghirlanda--Guerra identities}

The identity below originates in the auxiliary-interaction argument of
Ghirlanda and Guerra \cite[\S3, equation~(28)]{GhirlandaGuerra1998}.
Here its conditioning includes the joint collection of overlap components.

A Gram array is a symmetric array $(R_{lj})_{l,j\geq1}$ whose finite
principal submatrices are positive semidefinite. It is weakly exchangeable
if its law is invariant under finite permutations of its indices. For a
finite vector $\mathbf R=(R^1,\ldots,R^m)$ of such arrays, let
$\mathcal F_n$ be the sigma-field generated by all its entries with indices
at most $n$. We say that $\mathbf R$ satisfies the joint
Ghirlanda--Guerra identities if, for every $n\geq2$ and every bounded
measurable $f$,
\begin{equation}\label{arr:gg}
 \E[f(\mathbf R_{1,n+1})\mid\mathcal F_n]
 =\frac1n\E f(\mathbf R_{12})
   +\frac1n\sum_{l=2}^n f(\mathbf R_{1l}).
\end{equation}
The identities with any other distinguished index follow by
exchangeability. In all applications the arrays are bounded and each
diagonal is a deterministic constant.

The following invariance identity is the ingredient that turns
\eqref{arr:gg} into a geometric description of an array. It is the
array form of the invariance formula in
\cite[Theorem~2 in the cited arXiv version]{PanchenkoUltra}.

\begin{lemma}\label{arr:invariance}
Let $R$ be a bounded symmetric weakly exchangeable array satisfying
\eqref{arr:gg}. Fix $n\geq2$ and bounded measurable functions
$f_1,\ldots,f_n$. Define
\[
 H_l=\sum_{i=1}^n f_i(R_{il}),\qquad
 H'_j=H_j-f_j(R_{jj})+\E f_j(R_{12})\quad(1\leq j\leq n).
\]
For real $t$, let $Z(t)$ be the limiting empirical mean of $e^{tH_l}$
over $l>n$. Then, for every bounded $\mathcal F_n$-measurable $\Phi$,
\begin{equation}\label{arr:invariance-eq}
 \E\Phi
 =\E\left[
  \frac{\Phi\exp(t\sum_{j=1}^n H'_j)}{Z(t)^n}
 \right].
\end{equation}
\end{lemma}

\begin{proof}
An empirical mean of bounded exchangeable variables is Cauchy in $L^2$:
in its second-moment calculation all cross terms with distinct indices
coincide, and the contribution from equal indices tends to zero. Applying
this observation to the powers of $H_l$ constructs their empirical means
simultaneously. The uniformly convergent exponential series then defines
$Z(t)$ and agrees with the empirical mean of $e^{tH_l}$. In particular,
if $|H_l|\leq C$, then
$e^{-C|t|}\leq Z(t)\leq e^{C|t|}$.

Write $\varphi(t)$ for the right-hand side of
\eqref{arr:invariance-eq}. Differentiating a denominator can be represented
by averaging over one further replica. Iterating this operation gives
\[
 \varphi^{(d)}(t)
 =\E\left[
 \frac{\Phi\exp\bigl(t\sum_{l\leq n}H'_l
                   +t\sum_{n<l\leq n+d}H_l\bigr)}
      {Z(t)^{n+d}}
 \prod_{j=n}^{n+d-1}
 \left(\sum_{l\leq n}H'_l+
       \sum_{n<l\leq j}H_l-jH_{j+1}\right)
 \right].
\]
To justify the replacement by a new replica, average over indices beyond
all those already present. The empirical means are unchanged by finite
permutations of these indices, and boundedness permits passage to their
$L^2$ limits.

At $t=0$, condition on the first $n+d-1$ replicas. For each
$i\leq n$, Equation~\eqref{arr:gg} computes the conditional mean of
$f_i(R_{i,n+d})$. Summing those identities makes the last factor in the
product have conditional mean zero. Thus $\varphi^{(d)}(0)=0$ for
$d\geq1$.

This derivative argument determines the identity for all real $t$.
Indeed, near a fixed real $t_0$,
\[
 |Z(t_0+z)-Z(t_0)|
 \leq Z(t_0)(e^{C|z|}-1).
\]
Consequently the denominator is uniformly nonzero in a complex
neighborhood of $t_0$, and the expected ratio is analytic there by
boundedness. The vanishing derivatives prove constancy near zero;
overlapping such neighborhoods extends constancy along the real axis.
\end{proof}

The following duplication principle follows the invariance method of
\cite[Theorem~4 in the cited arXiv version]{PanchenkoUltra}.

\begin{lemma}[Duplication]\label{arr:duplication}
Let $R$ satisfy the hypotheses of Lemma~\ref{arr:invariance}.
Suppose an event $\mathcal A$ of positive probability prescribes open
tolerances for the off-diagonal entries among $n\geq2$ replicas.
Assume that its requirements on the pairs $(i,n)$, $i<n$, are intervals
$I_i\subset(-\infty,u)$. Then
\[
 \Prob\bigl(\mathcal A,\ R_{i,n+1}\in I_i\ (i<n),\
                         R_{n,n+1}<u\bigr)>0.
\]
\end{lemma}

\begin{proof}
Let $W$ be the empirical future frequency of
$\{R_{il}\in I_i\text{ for all }i<n\}$, and let $V$ be that of
$\{R_{nl}\geq u\}$. If the asserted probability were zero, then
$V\geq W$ almost surely on $\mathcal A$. Apply
Lemma~\ref{arr:invariance} with
$f_n=\ind_{[u,\infty)}$, all other $f_i=0$, and
$\Phi=\ind_{\mathcal A}$. Put $q=\Prob(R_{12}\geq u)$.
On $\mathcal A$ the numerator contributed by the $H'_j$ is $e^{tq}$.
For $t>0$, therefore,
\[
 \Prob(\mathcal A)
 \leq e^{tq}\E\frac{\ind_{\mathcal A}}
                         {(1+(e^t-1)W)^n}
 \leq e^{tq}\E\frac{W}{(1+(e^t-1)W)^n}.
\]
For the second inequality, discard the requirements among the first
$n-1$ replicas and exchange replica $n$ with a future index. The variable
$W$ only tests overlaps against the first $n-1$ replicas and is unchanged
by that exchange, so the empirical mean of the remaining indicator is
$W$.

The positive-probability event $\mathcal A$ implies $q<1$. For $n\geq2$,
\[
 \sup_{0\leq w\leq1}\frac{w}{(1+(e^t-1)w)^n}
 \leq \frac{(n-1)^{n-1}}{n^n(e^t-1)}.
\]
The preceding upper bound consequently tends to zero as $t\to\infty$,
contradicting $\Prob(\mathcal A)>0$.
\end{proof}

\begin{proposition}[Structure of joint overlap arrays]\label{arr:structure}
Let $\mathbf R=(R^1,\ldots,R^m)$ be bounded, jointly weakly exchangeable
Gram arrays with deterministic diagonals, satisfying
\eqref{arr:gg}. Each scalar component has nonnegative off-diagonal
entries and is ultrametric:
\[
 R^a_{23}\geq\min(R^a_{12},R^a_{13})\quad\hbox{almost surely}.
\]
The support of the law of $\mathbf R_{12}$ is ordered coordinatewise.
Hence this law has a representation by simultaneous nondecreasing
quantiles $p_a:(0,1)\to\R_+$. If $Q=\sum_a R^a$ and $p$ is its
off-diagonal quantile, then $p=\sum_a p_a$, and every $p_a$ is a
measurable function of $p$. In particular, each $p_a$ is constant almost
everywhere on every interval on which $p$ is constant.
\end{proposition}

\begin{proof}
We first prove the scalar assertions. If
$r=\Prob(R_{12}< -\epsilon)>0$, then, on an event where every distinct
overlap among $n$ replicas is below $-\epsilon$, Equation~\eqref{arr:gg}
gives
\[
 \Prob(R_{i,n+1}\geq-\epsilon\mid\mathcal F_n)
 =\frac{1-r}{n}\quad(1\leq i\leq n).
\]
A union bound leaves conditional probability at least $r$ that all new
overlaps are below $-\epsilon$. Starting with two replicas, we obtain
arbitrarily large such lists with positive probability. If the diagonal
is bounded by $D$, a Gram representation $R_{ij}=v_i\cdot v_j$ of one
list would give
\[
 0\leq\left|\sum_{i=1}^n v_i\right|^2
 \leq nD-\epsilon n(n-1),
\]
which is impossible for large $n$. Thus the off-diagonal entries are
nonnegative.

For ultrametricity, we use the barycenter contradiction of
\cite[Section~3]{PanchenkoUltra}. Suppose ultrametricity fails. By
permuting indices there is a support
triple with
\[
 R_{12}=a\geq b=R_{13}>c=R_{23}.
\]
Apply Lemma~\ref{arr:duplication} repeatedly to copies of replicas 1
and 2. For any prescribed finite size and arbitrarily small tolerances,
this produces two sets of replicas whose cross overlaps are close to
$a$, whose internal off-diagonal overlaps are at most $a+\epsilon$,
and whose overlaps with replica 3 are close respectively to $b$ and $c$.
Here is the tolerance choice underlying this iteration. At each step
choose a support point in the current positive-probability matching
event, and choose the duplication threshold slightly larger than all
overlaps of the replica being copied. Allocate a total increase less
than $\epsilon$ among the finitely many steps. The requirements against
the other set and replica 3 remain in their original small intervals.

Let $\bar v_1,\bar v_2$ be the barycenters of Gram vectors for sets of
size $M$. Their squared norms are at most
$a+\epsilon+D/M$, and their mutual inner product is at least
$a-\epsilon$. It follows that
$|\bar v_1-\bar v_2|^2\leq4\epsilon+2D/M$. Their inner products with
the vector for replica 3, however, differ by a number arbitrarily close
to $b-c>0$. First take $M$ large and then the tolerances small to
contradict Cauchy--Schwarz. This proves scalar ultrametricity.

For synchronization, we use the coordinate-and-sum argument of
\cite[Section~4, Lemma~2]{PanchenkoMulti}. Each component and each sum
of two components is again a Gram array satisfying the scalar identities.
Suppose two points in the support of
the joint pair law have strictly crossed coordinates $a$ and $b$.
The independent-law term in Equation~\eqref{arr:gg}, with $n=2$,
gives positive probability of observing these two points in arbitrarily
small neighborhoods on pairs 12 and 13. Ultrametricity in each component
forces the corresponding coordinate on pair 23 to equal the smaller of
the two coordinates. The sum of coordinates $a$ and $b$ on pair 23 is
then strictly below its values on both pairs 12 and 13, contrary to
ultrametricity for that sum. The support is therefore ordered.

An ordered compact support has at most one point with any prescribed
sum of its coordinates. Its coordinates are nondecreasing functions
of that sum; in fact each is 1-Lipschitz on the set of possible sums.
Applying these functions to the quantile of $Q_{12}$ supplies the
simultaneous quantiles and all the stated conclusions.
\end{proof}

\subsection{Finite cascades and hierarchical marks}

We use Ruelle's Poisson cascade construction \cite{Ruelle1987}.
The conditional marking argument follows the Poisson-rescaling calculus
of Bolthausen and Sznitman \cite[Appendix, Proposition~A.2]{BolthausenSznitman};
the proof below states the fractional-moment conditions needed for the
possibly unbounded marks used here.

Proposition~\ref{arr:structure} identifies the order of the overlaps.
We now construct the finite-step arrays and describe how their marks
change under a Gibbs weight. This calculation will apply both to
Gaussian linear fields and to the quadratic Gaussian factors in the
cavity argument.

Fix the interval weights $w_0,\ldots,w_{k-1}>0$ and endpoints
$0=\zeta_0<\zeta_1<\cdots<\zeta_k=1$. At every vertex at depth $i-1$,
place an independent Poisson point process $(v_{\alpha r})_{r\geq1}$
on $(0,\infty)$ with intensity
$\zeta_i v^{-1-\zeta_i}\,dv$, for $1\leq i\leq k-1$.
The mass of a leaf $\alpha\in\N^{k-1}$ is the product of the points on
its path. Let $T$ be the sum of these masses and let $\nu_\alpha$ be
the masses divided by $T$. When $k=1$, there is one leaf of mass one
and $T=1$. We refer to $(\nu_\alpha)$ as the finite Ruelle probability
cascade with these parameters.

For independently sampled leaves, let $L_{lj}$ be the depth of their
last common vertex, with $L_{lj}=k-1$ when their leaf labels agree.
A step overlap quantile with values $q_i$ on the intervals of lengths
$w_i$ is represented by $R_{lj}=q_{L_{lj}}$ for $l\ne j$;
Lemma~\ref{arr:rpc-gg} proves the required pair probabilities and GG
identities. The replica diagonal is prescribed separately. In particular,
two distinct replicas may choose the same leaf and have overlap
$q_{k-1}$ even when their self-overlap is larger.

We record why the normalization is legitimate. For $0<\zeta<1$,
the sum of a Poisson process of intensity $\zeta v^{-1-\zeta}dv$,
multiplied by independent positive marks $V$ with
$\E V^\zeta<\infty$, has Laplace transform
\[
 \exp\{-c_\zeta\E[V^\zeta]s^\zeta\},\qquad s\geq0,
\]
by direct integration of $1-e^{-svV}$. It is finite and strictly
positive almost surely. Integrating this Laplace transform, or its
complement, against powers of $s$ gives every negative moment and every
positive moment of order below $\zeta$. In particular its logarithm
has finite second moment. Starting at the deepest level and proceeding
upwards, the strict inequalities between the $\zeta_i$ give all required
fractional moments. Thus $0<T<\infty$ almost surely and
$\E|\log T|^2<\infty$.

The Poisson facts just used require no additional construction theorem:
on a set of finite intensity, draw a Poisson number of independent
points with normalized intensity. Summing its exponential series gives
the Laplace formula and the mapping and independent-mark rules.
Increasing finite-intensity restrictions proves the same rules for the
processes above.

\begin{proposition}[Marked cascade recursion]\label{arr:rpc}
Give the root a random mark $\eta_0$. Conditional on ancestor marks,
give each child at level $i$ a mark with probability kernel
$K_i(d\eta_i\mid\eta_0,\ldots,\eta_{i-1})$, independently for distinct
children, for $1\leq i\leq k-1$. All marks are independent of the
Poisson processes conditional on their ancestors. One may also integrate
a residual mark $\eta_k$ at each leaf, with kernel $K_k$, afresh on
every visit to that leaf.

For a measurable payoff $X(\eta_0,\ldots,\eta_k)$ define
\[
 X_{k-1}=\log\int e^X\,dK_k,
 \qquad
 X_{i-1}=\frac1{\zeta_i}\log\int e^{\zeta_i X_i}\,dK_i
 \quad(i=k-1,\ldots,1).
\]
Here $X=-\infty$ is allowed on excluded residual states, so the leaf
integral can impose a hard constraint without renormalizing its prior;
the integral itself must be strictly positive and finite. If there is
no residual mark, set $X_{k-1}=X$ instead. Assume that all
these quantities are finite real numbers almost surely under the
original path law, and that $\E|X_0|<\infty$. Then
\begin{equation}\label{arr:rpc-log}
 \E\log\sum_\alpha\nu_\alpha e^{X_{k-1}(\alpha)}=\E X_0.
\end{equation}
After tilting the leaf masses by $e^{X_{k-1}}$, the unmarked cascade
has its original law up to tree relabeling. Conditional on the
branching of sampled leaves, their marks use the kernels
\begin{equation}\label{arr:rpc-kernels}
 d\widetilde K_i
 =e^{\zeta_i(X_i-X_{i-1})}\,dK_i
 \quad(1\leq i\leq k-1),
 \qquad
 d\widetilde K_k=e^{X-X_{k-1}}\,dK_k.
\end{equation}
The root mark keeps its original law: the Gibbs probability is
normalized for each fixed root mark, and the outer expectation still
averages that mark under its original distribution. Different branches use
conditionally independent marks, and separate visits use independent
residual marks even when their leaf labels coincide. These assertions
hold for every bounded measurable test of the sampled branching and
marks.
\end{proposition}

\begin{proof}
Multiply each point on an edge to level $i$ by
$e^{X_i-X_{i-1}}$. Conditional on earlier marks,
\[
 \int e^{\zeta_i(X_i-X_{i-1})}\,dK_i=1.
\]
Poisson mapping therefore preserves the point intensity and changes
the conditional mark law to $\widetilde K_i$. Applying this observation
recursively, and ordering the transformed child points at each vertex,
gives an ordinary unmarked cascade independent of the root mark. The
transformed child marks have the kernels in
\eqref{arr:rpc-kernels}; transformed unmarked descendant processes are
independent of their ancestor marks. Positivity and finiteness of the
edge-kernel densities ensure that the almost-sure finiteness assumptions remain
valid under these kernels.

Let $T_{\rm new}$ denote the sum of the products of transformed edge
points. The multipliers telescope along every path, giving the almost
sure identity
\begin{equation}\label{arr:rpc-total}
 \log\sum_\alpha\nu_\alpha e^{X_{k-1}(\alpha)}
 =X_0+\log T_{\rm new}-\log T.
\end{equation}
Conditional on the root, both totals have the original law of $T$.
Thus their log difference has conditional mean zero and conditional
second moment at most $4\E|\log T|^2$. This proves
\eqref{arr:rpc-log}, including integrability. In particular no
unconditional first moment of $e^X$ is required.

The same mapping proves the sampled-mark assertion. Relabeling is
performed separately among children of each vertex, so it preserves
ancestors and common levels. The residual integration contributes
$e^{X_{k-1}}$ to a leaf mass; conditional residual sampling after the
tilt has density $e^{X-X_{k-1}}$. Independent integration on separate
visits gives precisely the last assertion.
\end{proof}

\begin{lemma}\label{arr:rpc-gg}
For the common levels $L_{lj}$ of independently sampled cascade leaves,
\[
 \Prob(L_{12}=i)=w_i\quad(0\leq i\leq k-1),
\]
and the common-level array satisfies the Ghirlanda--Guerra identities.
Any bounded symmetric weakly exchangeable scalar ultrametric array
satisfying those identities is
uniquely determined in law by its off-diagonal distribution and its
fixed diagonal. Consequently a nondecreasing finite-valued rounding of
its off-diagonal entries, with the fixed diagonal retained separately,
has the cascade law prescribed by its rounded quantile.
\end{lemma}

\begin{proof}
Attach independent standard Gaussian marks to the vertices at level
$i\geq1$, and tilt by $e^{cg(\alpha)}$, where $g(\alpha)$ is the mark
on the path to $\alpha$ and $c\ne0$. By
Proposition~\ref{arr:rpc}, the unmarked branching law is unchanged and
the marks are shifted by $c\zeta_i$. Their covariance before tilting is
$\ind_{\{L_{lj}\geq i\}}$. Gaussian integration by parts against a
bounded test $\Phi$ of the common levels of $n$ replicas gives
\begin{equation}\label{arr:rpc-gg-eq}
 \zeta_i\E\Phi
 =\E\left[\Phi\left(1+
    \sum_{j=2}^n\ind_{\{L_{1j}\geq i\}}
    -n\ind_{\{L_{1,n+1}\geq i\}}\right)\right].
\end{equation}
For this calculation one can first retain finitely many leaves; Gaussian
exponential moments and Jensen's inequality for inverse partition
factors justify passage to the full probability prior. At $n=1$,
Equation~\eqref{arr:rpc-gg-eq} gives
$\Prob(L_{12}\geq i)=1-\zeta_i$. For general $n$, it gives the
Ghirlanda--Guerra identity for each threshold, and hence for every
function of the finite set of levels.

For uniqueness, first consider an array with finitely many possible
off-diagonal values. At any one of its thresholds, ultrametricity makes
the relation of being above the threshold an equivalence relation on
the sampled indices. If its pair probability is $1-\zeta$, the
conditional probability that the next sample joins an existing class
$C$ is
\[
 \frac{1-\zeta+|C|-1}{n}=\frac{|C|-\zeta}{n},
\]
by applying the identities to a representative of $C$. The classes
at different thresholds are nested. These joining probabilities, and
their differences for nested classes, determine the distribution of
the whole next row of overlaps. Induction on the sample size gives
uniqueness. General bounded off-diagonal distributions follow by
nondecreasing finite-valued roundings and convergence on each finite
array. Equal rounded values may first be merged into a single level.
\end{proof}

\subsection{Gaussian marks and passage through finite replicas}

For a cascade with levels $0,\ldots,k-1$, let
$C_0,\ldots,C_{k-1},C_{\rm d}$ be symmetric matrices such that
\[
 0\leq C_0\leq C_1\leq\cdots\leq C_{k-1}\leq C_{\rm d}
\]
in positive semidefinite order. A centered Gaussian vector with
covariance $C_i$ between replicas of common level $i$, and diagonal
$C_{\rm d}$, is obtained by independent increments along the tree:
the common root increment has covariance $C_0$, and the increment at
level $i\geq1$ has covariance $C_i-C_{i-1}$. The remaining covariance
$C_{\rm d}-C_{k-1}$ is a residual integrated independently on every
sampled visit. This construction also covers singular covariance
matrices and tied levels.

\begin{lemma}[Finite-replica continuity]\label{arr:finite-replica}
For a sequence of random probability priors, draw conditionally
independent replicas and form their overlap arrays. Suppose the
finite-dimensional laws of these bounded overlap arrays and
their Gaussian covariance arrays converge jointly, and the covariance
diagonals are uniformly bounded. Conditional on these arrays, let the
additional marks be centered Gaussian vectors with those covariances.
Then their joint finite-replica laws converge. Moreover, the following
quantities pass to the limit whenever the factors and tests are bounded
continuous functions of the corresponding finite-replica data:
\begin{enumerate}
 \item the expected logarithm of a one-replica Gibbs integral of a
 factor $f$ satisfying $0<\delta\leq f\leq M<\infty$;
 \item any bounded finite-replica test under reweighting by such a
 factor.
\end{enumerate}
The same conclusion holds for finite-step cascade approximations whose
overlap and covariance arrays converge in this sense.
\end{lemma}

\begin{proof}
Finite-dimensional centered Gaussian laws depend continuously on their
covariance matrices, including at singular matrices. Conditional
integration followed by the assumed convergence proves the first
claim.

Write $Z=\langle f\rangle$ for the original Gibbs integral. It belongs
to $[\delta,M]$. Uniformly approximate $\log Z$, or $Z^{-n}$ for an
$n$-replica reweighted test, by a polynomial on this interval. Each
power of $Z$ is an integral over additional independent Gibbs replicas.
The Gaussian marks on these extra replicas are evaluations of the
\emph{same} field, so their correlations are exactly those of the
joint covariance array. Every resulting term is a bounded
finite-replica expectation and converges by the first claim. The
uniform polynomial errors can then be sent to zero.
\end{proof}

In particular, to approximate a quantile $p$ one may round its values
by nondecreasing finite-valued maps converging uniformly to the
identity. Lemma~\ref{arr:rpc-gg} identifies the rounded array with a
finite cascade. Gaussian covariance paths may be approximated at the
same time provided their increments and diagonal excess remain
positive semidefinite and the covariances converge. A covariance path
that is a function of $p$ must be rounded consistently on its
plateaus. Unbounded factors require additional truncation estimates;
Lemma~\ref{arr:finite-replica} does not assert convergence of their
untruncated moments.

\begin{lemma}[Gaussian differentiation]\label{arr:gaussian-calculus}
Let $\pi$ be a probability prior, possibly random. Conditional on the
prior and its other disorder, let $X_t(\sigma)$ be a centered Gaussian
field with differentiable covariance $C_t(\sigma,\tau)$, and let
$V_t(\sigma)$ be a differentiable deterministic exponent. Assume, on
each compact parameter interval, uniform bounds on
$C_t(\sigma,\sigma)$, $|\dot C_t(\sigma,\tau)|$, $|V_t(\sigma)|$,
and $|\dot V_t(\sigma)|$, and uniform continuity in $t$ of the
derivatives. Then
\begin{equation}\label{arr:gaussian-pressure}
 \frac{d}{dt}\E\log\int e^{X_t+V_t}\,d\pi
 =\frac12\E\langle\dot C_t(1,1)-\dot C_t(1,2)\rangle_t
   +\E\langle\dot V_t(1)\rangle_t.
\end{equation}
For a bounded test $\Phi$ of $n$ replicas with no explicit dependence
on these Gaussian variables or on $t$,
\begin{equation}\label{arr:gaussian-test}
 \left|\frac{d}{dt}\E\langle\Phi\rangle_t\right|
 \leq\|\Phi\|_\infty
 \left(2n(n+1)\|\dot C_t\|_\infty
                    +2n\|\dot V_t\|_\infty\right).
\end{equation}
These formulas apply to the finite cascades above, including at
one-sided parameter endpoints by integration and continuity.
\end{lemma}

\begin{proof}
For a finite prior, differentiate the Gibbs ratios and integrate by
parts in the Gaussian coefficients. The logarithm gives
\eqref{arr:gaussian-pressure}. The covariance part for an $n$-replica
test is one half the expectation of $\Phi$ times
\[
 \sum_{l,m\leq n}\dot C_t(l,m)
 -2n\sum_{l\leq n}\dot C_t(l,n+1)
 +n(n+1)\dot C_t(n+1,n+2)
 -n\dot C_t(n+1,n+1).
\]
The deterministic part inserts
$\sum_{l\leq n}\dot V_t(l)-n\dot V_t(n+1)$.
Taking absolute values proves the stated bound.

For a general probability prior, approximate its integral by finite
empirical samples, or by finite subsets whose prior mass tends to one.
At a fixed system size and on a bounded parameter interval, a uniform
variance bound for $X_t$ and a bound for $V_t$ give all required
polynomial-exponential moments. Inverse partition factors are controlled
by Jensen's inequality:
\[
 \left(\int e^{X_t+V_t}\,d\pi\right)^{-r}
 \leq\int e^{-r(X_t+V_t)}\,d\pi\qquad(r>0).
\]
H\"older's inequality gives the same control for products occurring in
the differentiated ratios. These bounds justify passage to the limit
in the integrated formulas, and then differentiation wherever the
covariance derivatives are continuous. The integrated identities also
give the asserted endpoint statements.
\end{proof}
